% Adaptación del preámbulo de Paper III v1.5 español; no plantilla standalone de Pandoc.
\documentclass[english,11pt,a4paper]{article}
\usepackage[margin=1in]{geometry}
\usepackage{fontspec}
\setmainfont{lmroman10-regular.otf}[BoldFont=lmroman10-bold.otf,ItalicFont=lmroman10-italic.otf,BoldItalicFont=lmroman10-bolditalic.otf]
\setsansfont{lmsans10-regular.otf}[BoldFont=lmsans10-bold.otf]
\setmonofont{lmmono10-regular.otf}[Scale=MatchLowercase]
\newlength{\LeanCharWidth}
\usepackage{amsmath,amssymb}
\usepackage[english]{babel}
\usepackage{microtype,parskip}
\usepackage{graphicx,float,caption}
\usepackage{longtable,booktabs,array,calc}
\usepackage{xcolor}
\usepackage{xurl}
\usepackage{hyperref,bookmark}
\usepackage{etoolbox,needspace}
\setcounter{secnumdepth}{-\maxdimen}
\setlength{\emergencystretch}{3em}
\providecommand{\tightlist}{\setlength{\itemsep}{0pt}\setlength{\parskip}{0pt}}
\newcounter{none}
\AtBeginEnvironment{longtable}{\Needspace{8\baselineskip}\small}
\pretocmd{\section}{\Needspace{5\baselineskip}}{}{}
\pretocmd{\subsection}{\Needspace{9\baselineskip}}{}{}
\pretocmd{\subsubsection}{\Needspace{6\baselineskip}}{}{}
\captionsetup{font=small,labelfont=bf}
\urlstyle{same}
\hypersetup{pdfauthor={Juan Pablo Traverso Gianini},pdflang={en-US},colorlinks=true,linkcolor=blue,urlcolor=blue,citecolor=blue}
\title{Paper VI\\Formalization and Reproducibility Companion}
\author{Juan Pablo Traverso Gianini}
\date{}\begin{document}\raggedbottom\maketitle
\textbf{Formalization companion}\\
\textbf{Preprint 1.0}\\
\textbf{Date:} October 18, 2026.

Copyright © 2026 Juan Pablo Traverso Gianini. Licensed under CC BY-NC 4.0.

\clubpenalty=10000\widowpenalty=10000
\subsection{1. Source package and entry
point}\label{source-package-and-entry-point}

The
\href{https://github.com/jtraverso/erdos-81-chordal-clique-partitions/tree/main/preprints/PAPER_VI/05_formalization}{public
source package} contains the formal statements and proofs for Paper VI.
Its default target, \texttt{PaperVIRelease}, publicly imports
\texttt{PaperVI}, the compatibility alias \texttt{PaperVIExtended}, and
\texttt{PaperVIRecourse}. The complete local import closure contains 79
modules. It includes Theorems A--E, the general-size window, intrinsic
and exceptional results, the appendices, and the complete-join all-K5
prefix, recourse and restricted-connectivity results of Section 10.

The toolchain is \path{leanprover/lean4:v4.35.0}; Mathlib is pinned to
\path{56a005855ea17e0b0958f73ec801ae140b0a1679}. The author's Lean
sources are licensed under Apache-2.0. Copyright and SPDX headers
preserve existing author credits and third-party notices. The theorem
types and proof bodies are unchanged by these licensing headers.

\path{PaperVIRelease.lean} SHA256:
\path{05de8cc26e4077d04e74d9a15c1373e40f8c66ed32590df90c617509bf44acee}.

\path{SOURCE_CLOSURE.json} SHA256:
\path{1c121b7b8bf561acde97166ee7fe7c16b68d7dba771d601a0705633b396227a3}.

\path{PAPER_VI_Preprint1.0_LEAN435_FORMAL_SOURCE_FREEZE.zip}
SHA256:
\path{c5c32a9d9315350d8e2683f0c58389824116bd1844f58f7b09ad78d7fe3b61bb}.

\subsection{2. Theorem correspondence}\label{theorem-correspondence}

The printed theorem letters follow the manuscript's reading order. The
Lean identifiers below are the stable names in the supplied source
archive.

{\def\LTcaptype{none} % do not increment counter
\begin{longtable}[]{@{}
  >{\raggedright\arraybackslash}p{(\linewidth - 2\tabcolsep) * \real{0.5000}}
  >{\raggedright\arraybackslash}p{(\linewidth - 2\tabcolsep) * \real{0.5000}}@{}}
\toprule\noalign{}
\begin{minipage}[b]{\linewidth}\raggedright
Manuscript result
\end{minipage} & \begin{minipage}[b]{\linewidth}\raggedright
Lean declaration
\end{minipage} \\
\midrule\noalign{}
\endhead
\bottomrule\noalign{}
\endlastfoot
A: exact local formula & \texttt{PaperVI.theoremA} \\
B: restriction to triangles & \texttt{PaperVI.theoremC} \\
C: exterior prescriptions & \texttt{PaperVI.theoremB} \\
Proposition 2.1: general-size window &
\path{PaperVI.cp_window_K4_simple}, with the order-four
construction in \path{PaperVI.cp_window_K4} \\
Second complete-split baseline &
\path{PaperVI.Research.second_split_baseline} \\
Perturbations of the second core &
\path{PaperVI.Research.second_core_formula} \\
Intrinsic maximum & \path{PaperVI.Research.cp_eq_max_PhiQ} \\
Degree recognition & \path{PaperVI.Research.core_eq_highDegree} \\
D: exceptional partition optimum &
\path{PaperVI.Research.ExSetup.theoremE1} \\
D: fractional equality optimum &
\path{PaperVI.Research.ExSetup.theoremE1_frac} \\
D: triangle optimum & \path{PaperVI.Research.ExSetup.theoremE2} \\
D: restriction gap & \path{PaperVI.Research.ExSetup.theoremE3} \\
Maximum over all cores in the exceptional family &
\path{PaperVI.Research.ExSetup.theoremE4} \\
Concrete linear-gap examples &
\path{PaperVI.Research.exModel_linear_gap} \\
Retaining a partial physical lifting &
\path{PaperVI.Research.theoremB_extending} \\
Exterior recourse & \path{PaperVI.Research.theoremB_recourse} \\
E: count for one prescription &
\path{PaperVI.Research.prescribed_count} \\
E: total lower and upper bounds & \path{PaperVI.Research.theoremN} \\
Lemma 10.1: exact compensation &
\path{PaperVIRecourse.used_add_resid_eq} \\
Lemma 10.1: arbitrary legal prefix &
\path{PaperVIRecourse.exists_optimal_partition_extending_prefix} \\
Physical compatible-triangle prefix &
\path{PaperVIRecourse.exists_optimal_partition_of_compatible_triangles} \\
Proposition 10.2: full 34k/68k recourse &
\path{PaperVIRecourse.exists_full_allK5_recourse} \\
Optimal terminal counts &
\path{PaperVIRecourse.optimal_allK5_counts} \\
Protected lifting repair &
\path{PaperVIRecourse.protected_allK5_step} \\
First connectivity phase & \path{PaperVIRecourse.lift_alignment} \\
Second phase: steps of distance 18 &
\path{PaperVIRecourse.terminal_alignment_missing_bound},
\path{PaperVIRecourse.terminal_alignment} \\
Theorem 10.3: restricted connectivity &
\path{PaperVIRecourse.restricted_allK5_connectivity} \\
Balanced reference identity and unperturbed split optimum &
\texttt{PaperVI.Mval\_eq}, \path{PaperVI.cp_refGraph_eq_Mval} \\
Universal bipartition certificate and lower bound &
\path{PaperVI.slack_identity}, \path{PaperVI.slack_nonneg},
\texttt{PaperVI.cp\_ge} \\
Triangle packing: integral and real fractional optima &
\path{PaperVI.triangle_packing_eq} \\
Lemma 6.1: rigidity of the chordal core &
\path{PaperVI.core_complete_of_chordal} \\
Chordal formula and two-sided edit stability &
\texttt{PaperVI.cp\_chordal}, \path{PaperVI.chordal_stability} \\
Real fractional partition optimum & \path{PaperVI.fractional_eq_cp},
\path{PaperVI.frac_slack_identity} \\
Integral and fractional gain optima & \texttt{PaperVI.gain\_eq} \\
Types in every integral and fractional optimum &
\path{PaperVI.optimal_types},
\path{PaperVI.frac_optimal_support} \\
Exceptional partial-star obstruction &
\path{PaperVI.Research.ExSetup.partial_star_obstruction} \\
Appendix A: elementary complete-join formula at p\ensuremath{\ge}24Δ(H) &
\texttt{PaperVI.theoremA1} \\
Appendix B: localization of inefficient blocks on at most 21t edges &
\path{PaperVI.inefficient_edges_le}, using
\path{PaperVI.choose_two_le_21_slack} \\
Distinct optimal prescriptions &
\path{PaperVI.two_pow_m_optimal} \\
\end{longtable}
}

The companion map covers 39 manuscript result groups.
\path{COMPLETE_THEOREM_MAP.json} identifies their source modules and
declaration locations; \path{SOURCE_IMPORT_DECLARATION_MAP.json}
covers every module in the imported closure. The local window in A and B
is p\ensuremath{\ge}18h+1 (p\ensuremath{\ge}18h in the balanced case); C uses p\ensuremath{\ge}63h. Proposition 2.1
has its separately stated general-size inequalities.

All partition conclusions concern actual clique partitions of G. The
exceptional setup specifies a clique P, an independent R, complete
adjacency between them, and one additional vertex z with the specified
neighborhood. It is not merely an integer model with an assumed
partition number.

\subsection{3. Exact scope of the encoded
conclusions}\label{exact-scope-of-the-encoded-conclusions}

Theorem D and the partial-star example in Section 8.3 require
\texttt{\textbar{}R\textbar{}\ \ensuremath{\ge}\ \textbar{}P\textbar{}+2}, equivalently
\texttt{q\ \ensuremath{\ge}\ p+3}; its unrestricted-size setup permits p=0 and empty
neighbor sets. The maximum-over-cores comparison requires \texttt{q=2p}
and \texttt{p\ensuremath{\ge}19}. The concrete linear-gap family takes \texttt{p=3t},
\texttt{t\ensuremath{\ge}7}.

Recognition by degrees under A assumes n\ensuremath{\ge}3. It identifies the reference
satisfying the local theorem, but does not assert that the intrinsic
maximum has a unique maximizer. No uniqueness theorem for the intrinsic
maximizer is asserted.

Proposition 9.1 is conditional on the actual lifting invariant: legal
edges, disjoint used core pairs, disjoint core sets for incident
exterior edges, prescribed cardinalities and bounded load. In its
general mixed setting the terminal completion may change globally.
Thirteen counts exterior assignments, not whole-partition blocks.

Section 10 assumes a complete join and all exterior blocks of order
five. Proposition 10.2 assumes two already legal complete liftings and
q\ensuremath{\ge}2p−1; it has no small-load hypothesis or upper bound on q. Lean
carries arbitrary finite load bounds for these liftings, without
requiring a uniform bound in terms of h. Its retained-family conclusion
is physical: old terminal triangles avoiding the new lifting are
retained. The 34k/68k bounds apply to the full families of vertex-set
blocks, not merely an abstract coloring distance.

Theorem 10.3 additionally requires 2p−1\ensuremath{\le}q\ensuremath{\le}2p+1, p\ensuremath{\ge}63h, and load\ensuremath{\le}7h in
both endpoint states and throughout the path. In a complete join
h=editDegree=Δ(H). \texttt{BlockPath} records actual partitions at every
state and block symmetric difference at every step. The induction first
protects target-aligned exterior assignments, then protects shared
target terminal triangles. This is restricted connectivity; mixed K4/K5
optima and unbounded-load all-K5 optima remain outside its conclusion.
The stronger terminal-phase endpoint bounds the number of steps by the
number of target terminal triangles missing from the initial state. No
sharpness or running-time claim is made for these recourse and path
constants.

Theorem E counts labeled families of vertex sets, not ordered lists or
isomorphism classes. The formal statement gives the explicit finite
bounds, including the floor in the lower exponent. The elementary
logarithmic asymptotic is deduced in the manuscript; it is not a
separate formal premise.

Fractional partitions have equality on every edge and real nonnegative
weights. This is not fractional covering with inequalities and not
integrality of the entire polytope for arbitrary objectives. The
triangle-packing fractional optimum is likewise over the canonical real
model.

\subsection{4. Proved dependencies, build and axiom
audit}\label{proved-dependencies-build-and-axiom-audit}

Häggkvist--Janssen and ordinary Vizing are proved dependencies in the
supplied import closure, not assumed mathematical axioms. The former
supplies the residual list-edge-coloring completion. The latter supplies
the counting step. The author's registered formalizations are
Häggkvist--Janssen, Palomar \path{PALOMAR-2026-10-15-000003}, version
1, repository \path{jtraverso/haggkvist-janssen}, commit
\path{807e08bd7c480d4fae80c1b78377d23989d39ab0}; and Vizing, Palomar
\path{PALOMAR-2026-10-15-000004}, version 1, repository
\texttt{jtraverso/vizing}, commit
\path{01c72ed7efde2353a4bb1654f69b5c0bc5038863}. Registration is
provenance; the final local build and transitive axiom audit establish
the status of the dependencies actually imported here. No equitable
list-completion theorem is inferred from ordinary Vizing.

The aggregate build for Preprint 1.0 and the separate exported-type
check both exit zero. The check covers 32 exported types and their axiom
footprints. The all-declaration audit checks 1836 local declarations in
79 modules and visits 3421 declarations transitively, with zero
nonstandard axioms. Its axiom union is exactly \texttt{propext},
\texttt{Classical.choice} and \texttt{Quot.sound}. Every listed module
must be present in the imported environment; a missing module causes the
audit to fail. A separate map matches source files, actual imports and
declaration counts. An active-source scan, masking comments and strings,
finds no \texttt{sorry}, \texttt{admit}, project \texttt{axiom},
\texttt{native\_decide}, \texttt{implemented\_by} or \texttt{extern}
escape.

The build used the pinned Mathlib cache and compatible local build
outputs. Lake rebuilt the sources with new licensing headers and their
affected imports. The final checks ran against that aggregate. It was
not a cache-free reconstruction. Deprecation and inherited linter
warnings remain in the logs; no warning-free claim is made.
Authoritative logs are \path{BUILD_PUBLIC_1_0.portable.log},
\path{ENDPOINTS_PUBLIC_1_0.portable.log} and
\path{FULL_AUDIT_PUBLIC_1_0.portable.log}, each with its recorded
zero exit. The source manifest and freeze receipt bind these records to
the final sources.

\subsection{5. Reproducing the formal
checks}\label{reproducing-the-formal-checks}

Extract
\path{formal_sources/PAPER_VI_Preprint1.0_LEAN435_FORMAL_SOURCE_FREEZE.zip}
into a fresh directory, install its pinned Lean toolchain, and run the
following from that project root. The archive contains all local sources
and configuration, but no \texttt{.lake} cache or machine-specific
junction.

\begin{verbatim}
lake exe cache get
lake build PaperVIRelease
lake env lean -j1 -s32768 EndpointGate.lean
lake env lean -j1 -s32768 FreezeFullAudit.lean
\end{verbatim}

Keep the supplied \texttt{lake-manifest.json} unchanged. The cache
command obtains compatible files for those pinned dependencies; do not
delete or update the manifest to select dependency versions. A build of
a single intermediate module is not a substitute for the canonical
aggregate and its two final checks.

\subsection{6. Finite certificates and adversarial regression
checks}\label{finite-certificates-and-adversarial-regression-checks}

The finite-certificate package in \path{03_reproducibility/} contains
33 literal-graph certificates: 18 local A--C cases and 15 exceptional
cases. Local cases include the balanced boundary p = 18h, the general
boundary p = 18h + 1, all three allowed exterior sizes, simultaneous
core and crossing deletions with exterior additions, concentrated
degree-two stars, an exterior cycle, and mixed K4/K5 prescriptions at p
= 63h. The exceptional cases use p = 2, q = 5, all a = 0,1,2 and b =
0,1,2,3,4; they include empty neighbor sets and the a = 2, b = 1 K4
branch.

The generator uses a finite greedy completion heuristic, not an
implementation of the universal Häggkvist--Janssen proof. An independent
checker imports no generator code. It reconstructs every clique and edge
multiplicity, checks types, prescriptions and exact counts, and checks
signed dual bounds for the exceptional fractional partition objective.
It also checks actual triangle partitions against nonnegative
triangle-cover duals. These finite primal--dual equalities certify the
examples; the universal results have their separate written and Lean
proofs.

Certo generates and replays six exact-cover certificates against their
exact specifications: a three-edit-class A case, a concentrated-star
triangle case, two boundary prescription cases, and two exceptional
cases. Certo's exact-cover result verifies the supplied partition;
global optimality in those finite cases is justified separately by the
checked integer dual. Three independent corruptions are rejected:
missing coverage, repeated coverage and an incorrect prescription. A
fourth Certo control changes the specification identity and is refused
by provenance checking. That refusal is not reported as a solver
rejection or a mathematical counterexample.

A separate two-exterior-edge complete-join example has three actual
optimal states. One K5 lifting is protected while the other changes. The
checker confirms 1938 retained compatible terminal triangles, 24 lost
old blocks, symmetric difference 48, and a final terminal recoloring of
distance 6. All states lie in the stated bounded-load all-K5 class; the
two-step path fits the finite bound 1949 and the uniform step bound 884.
Understated distance and load bounds are rejected as negative controls.
Three additional Certo certificates validate exact edge coverage in each
supplied state. Path distances, preservation and bounded-load membership
are checked independently. This is a finite path example, not an
exhaustive connectivity computation.

Run the following from \path{03_reproducibility/}. The independent
Python checkers use only the standard library. The Certo replay
additionally requires the recorded Certo installation. The supplied
certificates retain their specification identity; regeneration creates
certificates for the local specification paths.

\begin{verbatim}
python generate_v06.py
python verify_v06.py
python recourse_v06.py
python verify_recourse_v06.py
python run_certo_v06.py
python run_certo_recourse_v06.py
\end{verbatim}

No finite census is a proof premise of a universal theorem. The supplied
scripts and independently checked certificates support only their stated
finite instances.

\subsection{7. Mathematical and bibliographic
scope}\label{mathematical-and-bibliographic-scope}

The source package proves the encoded statements with their explicit
hypotheses. It does not establish a historical-priority claim or replace
independent mathematical review. The directed comparison includes
full-text checks of EFO, CEO, the EOZ author galley, Ma's threshold
bounds, McGuinness--Rees's line-graph enumeration and Wallis--Zhang's
independent-join characterization. Wallis--Zhang's Proposition 3.2 and
Corollary 3.3 include the boundary of Theorem D in which the exceptional
vertex is adjacent to every ordinary exterior vertex.

The relevant Pullman--Donald result is credited through the inspected
EFO article (p.~94) and Wallis--Zhang (Corollary 3.3); its original
article was not available for direct inspection. Wallis--Wu91 is
compared through its primary abstract and the authors' announcement.
These limitations preclude an exhaustive historical-priority assertion.
Okechukwu's Lemma 4.2 supplies a triangle-partition upper bound under
its own coloring, capacity and defect hypotheses; overlap with Theorem B
is asserted only in the subcases satisfying those hypotheses, not across
the full local-edit window.

Paper IV proves the sharp bound for sufficiently large chordal graphs,
together with an all-orders bound with an absolute additive constant.
Paper VI studies exact formulas and additional structure under its own
hypotheses. The separate IV--VI application note has its own entry point
and validation; this source package does not certify that note's source
closure. Arbitrary mixed K4/K5 connectivity and unbounded-load
connectivity are outside Section 10.

\end{document}
